\section{Arhitektura sustava}
\label{sec:02-arhitektura}

% izvor: launch/mission.launch.py:1-107; config/bridge.yaml; config/controllers.yaml; README.md §8
Arhitektura programskog sustava robota PAS-DUAL-ARM strukturirana je u pet jasno odvojenih
funkcionalnih slojeva (slika~\ref{fig:arh-sustava}): simulacijski sloj u Gazebu, sloj mostova i
hardverske apstrakcije, sloj niskorazinskih kontrolera (\texttt{ros2\_control}), sloj autonomnih
stogova za planiranje (MoveIt~2 i Nav2) te vršni aplikacijski sloj misijskih čvorova. Ovakva
razdioba osigurava modularnost, jednostavnu zamjenu simuliranih senzora i aktuatora stvarnim
sklopovljem te strogu izolaciju odgovornosti pojedinih podsustava.

\begin{figure}[H]
  \centering
  % G04: Arhitektura sustava (slojevi Gazebo, bridge/ros2_control, MoveIt/Nav2, aplikacijski cvorovi)
% Izvori: launch/mission.launch.py, launch/sim.launch.py, launch/nav2.launch.py, launch/task.launch.py,
%         config/bridge.yaml, config/controllers.yaml, pas_dual_arm_scripts
\begin{tikzpicture}[
  scale=0.97, transform shape,
  font=\scriptsize,
  layer/.style={draw, rounded corners=3pt, fill=#1, inner sep=4pt},
  nodebox/.style={draw, rounded corners=2pt, fill=white, align=center, inner sep=2pt, font=\tiny},
  arr/.style={-{Stealth[length=1.5mm]}, semithick},
  darr/.style={{Stealth[length=1.5mm]}-{Stealth[length=1.5mm]}, semithick}
]

  % Sloj 5: Aplikacija i misija (vrh)
  \begin{scope}[local bounding box=applayer]
    \node[nodebox, fill=blue!10, text width=24mm] (mt) at (1.5, 0) {\textbf{main\_task}\\orkestracija misije};
    \node[nodebox, fill=blue!10, text width=22mm] (rn) at (4.4, 0) {\textbf{room\_navigator}\\graf soba, vrata};
    \node[nodebox, fill=blue!10, text width=22mm] (nz) at (7.1, 0) {\textbf{nav\_zones}\\potencijalno polje};
    \node[nodebox, fill=blue!10, text width=22mm] (ad) at (9.8, 0) {\textbf{aruco\_detector}\\detekcija 4 markera};
    \node[nodebox, fill=blue!10, text width=22mm] (cm) at (12.5, 0) {\textbf{cmd\_vel\_relay}\\i collision\_monitor};
    \node[nodebox, fill=blue!5, text width=15mm] (gui) at (14.8, 0) {\textbf{nav\_gui}\\upravljač};
  \end{scope}

  % Sloj 4: Autonomni stogovi (MoveIt 2 & Nav2)
  \begin{scope}[local bounding box=stacklayer]
    \node[nodebox, fill=orange!15, text width=38mm] (moveit) at (3.0, -1.35)
      {\textbf{MoveIt 2} (\texttt{move\_group})\\kinematika KDL, OMPL, planiranje};
    \node[nodebox, fill=orange!15, text width=44mm] (nav2) at (8.5, -1.35)
      {\textbf{Nav2 Stog}\\NavfnPlanner, MPPI/DWB, costmapovi};
    \node[nodebox, fill=orange!15, text width=32mm] (amcl) at (13.5, -1.35)
      {\textbf{AMCL / slam\_toolbox}\\lokalizacija na karti i SLAM};
  \end{scope}

  % Sloj 3: ros2_control & Hardware Abstraction
  \begin{scope}[local bounding box=controllayer]
    \node[nodebox, fill=green!15, text width=42mm] (armctrl) at (3.0, -2.65)
      {\textbf{Ruke, vodilice, glava i hvataljke}\\JTC (pozicijski/effort) i GripperAction};
    \node[nodebox, fill=green!15, text width=44mm] (basectrl) at (8.5, -2.65)
      {\textbf{base\_controller}\\MecanumDriveController (\texttt{velocity})};
    \node[nodebox, fill=green!15, text width=32mm] (cmgr) at (13.5, -2.65)
      {\textbf{controller\_manager}\\100~Hz, \texttt{joint\_states}};
  \end{scope}

  % Sloj 2: Mostovi prema simulatoru
  \begin{scope}[local bounding box=bridgelayer]
    \node[nodebox, fill=purple!12, text width=70mm] (ignctrl) at (4.5, -3.85)
      {\textbf{ign\_ros2\_control} (\texttt{IgnitionSystem})\\dvosmjeran prijenos naredbi i stanja zglobova};
    \node[nodebox, fill=purple!12, text width=75mm] (gzbridge) at (11.5, -3.85)
      {\textbf{ros\_gz\_bridge} (\texttt{bridge.yaml})\\lidar, RGB-D kamere, kontakti, FT, attach/detach, clock};
  \end{scope}

  % Sloj 1: Gazebo simulator (dno)
  \begin{scope}[local bounding box=gzlayer]
    \node[nodebox, fill=gray!18, text width=152mm] (gzsim) at (7.8, -4.95)
      {\textbf{Gazebo Fortress (Ignition 6)} --- fizikalni pogon DART\\
       Svijet \texttt{seminar\_world.sdf}, senzorski dodaci (\texttt{gpu\_lidar}, kamere, kontakti, FT), \texttt{DetachableJoint}};
  \end{scope}

  % Oznake slojeva lijevo
  \node[font=\tiny\bfseries, text=blue!70!black, rotate=90] at (-0.3, 0) {APLIKACIJA};
  \node[font=\tiny\bfseries, text=orange!70!black, rotate=90] at (-0.3, -1.35) {STOGOVI};
  \node[font=\tiny\bfseries, text=green!50!black, rotate=90] at (-0.3, -2.65) {UPRAVLJANJE};
  \node[font=\tiny\bfseries, text=purple!70!black, rotate=90] at (-0.3, -3.85) {MOST};
  \node[font=\tiny\bfseries, text=gray!80!black, rotate=90] at (-0.3, -4.95) {SIMULACIJA};

  % Poveznice
  \draw[darr] (mt) -- (moveit);
  \draw[darr] (rn) -- (nav2);
  \draw[arr] (nz) -- (nav2);
  \draw[arr] (ad) -- (mt);
  \draw[arr] (cm) -- (basectrl);
  \draw[arr] (amcl) -- (nav2);
  \draw[darr] (moveit) -- (armctrl);
  \draw[darr] (nav2) -- (cm);
  \draw[darr] (armctrl) -- (ignctrl);
  \draw[darr] (basectrl) -- (ignctrl);
  \draw[darr] (cmgr) -- (ignctrl);
  \draw[darr] (ignctrl) -- (gzsim);
  \draw[arr] (gzsim) -- (gzbridge);
  \draw[arr] (gzbridge) -- (ad);
  \draw[arr] (gzbridge) -- (amcl);
  \draw[arr] (gzbridge) -- (cm);
\end{tikzpicture}

  \caption{Nacrt arhitekture sustava: pet slojeva od Gazebo fizike do aplikacijske orkestracije}
  \label{fig:arh-sustava}
\end{figure}

\subsection{Slojevi sustava i raspodjela čvorova}
\label{sec:02-slojevi}

% izvor: config/bridge.yaml:1-164; robot.urdf.xacro:479-481; launch/sim.launch.py:227-252
Simulacijski sloj izvodi se unutar simulatora Gazebo Fortress uz DART fizikalni rješavač.
Između Gazeba i ROS~2 ekosustava djeluju dva komplementarna mosta. Dodatak \texttt{ign\_ros2\_control}
\cite{mail2026} preko sučelja \texttt{IgnitionSystem} povezuje \texttt{controller\_manager}
s virtualnim zglobovima robota na frekvenciji od 100~Hz. Svi senzorski tokovi (lidar, RGB-D kamere,
kontaktni senzori na prstima, dvoosni tenzometrijski FT senzori) te naredbe za kruto spajanje
kocke (\texttt{DetachableJoint}) prenose se putem paketa \texttt{ros\_gz\_bridge} prema
deklaraciji u \texttt{bridge.yaml}.

% izvor: pas_dual_arm_scripts/setup.py:28-44; main_task.py; room_navigator.py; nav_zones.py
U tablici~\ref{tab:cvorovi} prikazani su ključni aplikacijski čvorovi razvijeni u sklopu paketa
\texttt{pas\_dual\_arm\_scripts}, njihove primarne odgovornosti te ulazno-izlazna komunikacijska sučelja.

\begin{table}[H]
  \centering
  \scriptsize
  \setstretch{1.05}
  \caption{Pregled aplikacijskih čvorova sustava u paketu \texttt{pas\_dual\_arm\_scripts}}
  \label{tab:cvorovi}
  \begin{tabularx}{\textwidth}{@{}>{\bfseries}p{2.8cm}>{\raggedright\arraybackslash}X>{\raggedright\arraybackslash}p{4.6cm}@{}}
    \toprule
    Čvor & Glavna odgovornost & Ključne teme, akcije i servisi \\
    \midrule
    \texttt{main\_task} & Orkestracija misije (8 koraka), dvoručni hvat, vizualno navođenje, odlaganje & Ulaz: \texttt{/mission/start}, detekcije, TF\newline Izlaz: \texttt{/room\_navigator/goto}, akcije ruku \\
    \texttt{room\_navigator} & Navigacija među sobama, prolazak kroz vrata uz provjeru poze ruku & Ulaz: \texttt{/room\_navigator/goto}, \texttt{/joint\_states}\newline Izlaz: \texttt{/navigate\_through\_poses}, status \\
    \texttt{nav\_zones} & Izračun jedinstvenog potencijalnog polja i brazdi nulte cijene iz karte & Ulaz: \texttt{/map}\newline Izlaz: \texttt{/nav\_zones/field} (costmap filter) \\
    \texttt{feature\_registry} & Semantičko prepoznavanje vrata, stolova i dock poza iz geometrije & Ulaz: \texttt{/map}\newline Izlaz: registar značajki, servisi koordinata \\
    \texttt{aruco\_detector} & Detekcija ArUco markera (glava i zapešća), estimacija 6D poze & Ulaz: \texttt{/camera/*}, \texttt{/wrist\_*/*}\newline Izlaz: \texttt{/aruco\_detections}, objave u \texttt{/tf} \\
    \texttt{cmd\_vel\_relay} & Arbitraža i sigurnosno prespajanje brzinskih naredbi baze & Ulaz: \texttt{/cmd\_vel\_safe}, \texttt{/cmd\_vel}\newline Izlaz: \texttt{reference\_unstamped} \\
    \texttt{collision\_monitor} & Neposredno zaustavljanje baze na temelju laserskog obrisa & Ulaz: \texttt{/scan\_filtered}, poligon obrisa\newline Izlaz: \texttt{/cmd\_vel\_safe} \\
    \texttt{footprint\_publisher} & Dinamičko prilagođavanje tlocrtnog poligona robota i nošene kocke & Ulaz: stanje misije, \texttt{/joint\_states}\newline Izlaz: \texttt{/footprint}, \texttt{/mission/carried\_points} \\
    \texttt{scan\_filter} & Uklanjanje točaka vlastitog tijela i kotača iz laserskog skena & Ulaz: \texttt{/scan}\newline Izlaz: \texttt{/scan\_filtered} \\
    \texttt{nav\_gui} & Tkinter grafičko sučelje za praćenje misije i ručni start & Ulaz: \texttt{/mission/task\_status}\newline Izlaz: objava na \texttt{/mission/start} \\
    \bottomrule
  \end{tabularx}
\end{table}

\subsection{Koordinatni sustavi i stablo transformacija (TF)}
\label{sec:02-tf}

% izvor: robot.urdf.xacro:20-250; kortex .../gen3.xacro; base.urdf.xacro; pan_tilt.xacro
Prostorni odnosi i kinematičke veze formalizirani su kroz stablo transformacija okvira (\texttt{/tf}
i \texttt{/tf\_static}). Ishodište navigacijskog prostora čini globalni okvir \texttt{map}, u kojem
algoritam AMCL procjenjuje pozu odometrijskog okvira \texttt{odom}. Kontroler baze
\texttt{mecanum\_drive\_controller} putem releja \texttt{cmd\_vel\_relay} objavljuje transformaciju
između \texttt{odom} i tlocrtne projekcije robota \texttt{base\_footprint}, na koju se nadovezuje
strukturna točka \texttt{base\_link}.

Kinematičko stablo robota grana se iz okvira \texttt{base\_link} na četiri glavne grane:
\begin{enumerate}
  \item \textbf{Lidar}: \texttt{base\_link} $\rightarrow$ \texttt{laser\_link} (postavljen na visini $0{,}209$~m, usmjeren vodoravno prema naprijed).
  \item \textbf{Torzo i vodilice}: \texttt{base\_link} $\rightarrow$ \texttt{torso\_link} $\rightarrow$ prizmatični zglobovi $\rightarrow$ klizači \texttt{torso\_left\_carriage\_link} i \texttt{torso\_right\_carriage\_link}.
  \item \textbf{Robotske ruke i hvataljke}: sa svakog klizača grana se kinematički lanac Kinova Gen3 manipulatora (\texttt{*\_base\_link} $\rightarrow$ zglobovi $1\dots7$ $\rightarrow$ \texttt{*\_end\_effector\_link} $\rightarrow$ \texttt{*\_bracelet\_link}), na koji su montirane kamere zapešća (\texttt{wrist\_*\_camera\_link}) te baze hvataljki Robotiq 2F-85 s kinematikom prstiju.
  \item \textbf{Pan-tilt glava}: s vrha torza (\texttt{torso\_link}) lanac vodi preko rotacijskih zglobova \texttt{pan\_tilt\_yaw\_link} i \texttt{pan\_tilt\_pitch\_link} do optičkog okvira RGB-D kamere \texttt{camera\_color\_optical\_frame}.
\end{enumerate}
Detektirani ArUco markeri i manipulirana kocka dinamički se vežu u TF stablo preko okvira kamere koja
ih trenutno opaža (\texttt{camera\_link} ili \texttt{wrist\_*\_link}), dok čvor \texttt{main\_task}
prilikom hvata pamti fiksnu transformaciju kocke u odnosu na zapešće (\texttt{\_hold}) kako bi
očuvao praćenje objekta kada marker izađe iz vidnog polja.

\subsection{Hijerarhija pokretačkih skripti (Launch)}
\label{sec:02-launch}

% izvor: launch/mission.launch.py:1-107; launch/sim.launch.py; launch/nav2.launch.py; launch/task.launch.py
Pokretanje cjelokupnog sustava izvedeno je hijerarhijski kroz četiri ugniježđene pokretačke skripte
(slika~\ref{fig:launch-hijerarhija}). Krovna datoteka \texttt{mission.launch.py} uvodi precizno
odmjerena vremenska kašnjenja kako bi se spriječilo zagušenje procesora i osigurala dostupnost
ROS~2 servisa prije pokretanja ovisnih čvorova:

\begin{figure}[H]
  \centering
  % G05: Hijerarhija pokretackih datoteka (mission.launch.py -> sim, nav2, task)
% Izvori: launch/mission.launch.py:20-107, launch/sim.launch.py, launch/nav2.launch.py, launch/task.launch.py
\begin{tikzpicture}[
  font=\scriptsize,
  topl/.style={draw, rounded corners=3pt, fill=blue!12, align=center, text width=55mm, inner sep=3pt},
  subl/.style={draw, rounded corners=2pt, fill=orange!15, align=center, text width=42mm, inner sep=3pt},
  nodebox/.style={draw, rounded corners=2pt, fill=white, align=center, text width=38mm, inner sep=2pt, font=\tiny},
  arr/.style={-{Stealth[length=1.5mm]}, semithick},
  lbl/.style={font=\tiny\bfseries, text=black!70, fill=white, inner sep=1pt}
]

  % Vrh: mission.launch.py
  \node[topl] (mission) at (7.5, 0)
    {\textbf{\texttt{mission.launch.py}}\\Krovna orkestracija pokretanja sustava};

  % Podređeni launch fajlovi
  \node[subl] (sim) at (2.2, -1.5)
    {\textbf{\texttt{sim.launch.py}} ($t = 0$~s)\\Simulacija i aktuatori};
  \node[subl] (nav) at (7.5, -1.5)
    {\textbf{\texttt{nav2.launch.py}} ($t = 6$~s)\\Mapiranje i navigacija};
  \node[subl] (task) at (12.8, -1.5)
    {\textbf{\texttt{task.launch.py}} ($t = 10$~s)\\Manipulacija i misija};

  % RViz (t = 8 s)
  \node[nodebox, fill=purple!10, text width=26mm] (rviz) at (10.15, -0.75)
    {\textbf{\texttt{rviz2}} ($t = 8$~s)\\prikaz \texttt{nav2.rviz}};

  % Čvorovi pod sim
  \node[nodebox] (gz) at (2.2, -2.65)
    {Gazebo Fortress (\texttt{seminar\_world.sdf})\\DART fizikalni rješavač};
  \node[nodebox] (spwn) at (2.2, -3.45)
    {Spawn robota iz \texttt{robot.urdf.xacro}\\početna poza AMCL-a $(0,0,0)$};
  \node[nodebox] (bridge) at (2.2, -4.25)
    {\texttt{ros\_gz\_bridge} (\texttt{bridge.yaml})\\senzori, kontakti, clock, spojevi};
  \node[nodebox] (ctrl) at (2.2, -5.05)
    {\texttt{spawner} (8 kontrolera)\\120~s timeout, \texttt{controllers.yaml}};

  % Čvorovi pod nav2
  \node[nodebox] (amclnode) at (7.5, -2.65)
    {Karta \texttt{seminar\_map.yaml}\\AMCL lokalizacija};
  \node[nodebox] (navstack) at (7.5, -3.45)
    {Nav2 navigacijski poslužitelji\\(planer, kontroler, costmapovi)};
  \node[nodebox] (navhelper) at (7.5, -4.25)
    {\texttt{feature\_registry}, \texttt{nav\_zones},\\i \texttt{room\_navigator}};
  \node[nodebox] (navsafety) at (7.5, -5.05)
    {\texttt{collision\_monitor}, \texttt{cmd\_vel\_relay},\\i \texttt{nav\_gui} (operater)};

  % Čvorovi pod task
  \node[nodebox] (mg) at (12.8, -2.65)
    {\texttt{move\_group} (MoveIt~2)\\planiranje dvoručnih poza};
  \node[nodebox] (aruco) at (12.8, -3.45)
    {\texttt{aruco.launch.py}\\4 instance \texttt{aruco\_detector}};
  \node[nodebox] (mtask) at (12.8, -4.45)
    {\textbf{\texttt{main\_task.py}} ($t = 13$~s)\\parametar \texttt{mission:=true}\\izvršavanje misijskog stabla};

  % Poveznice
  \draw[arr] (mission) -| node[lbl, pos=0.4] {odgoda 0~s} (sim);
  \draw[arr] (mission) -- node[lbl] {odgoda 6~s} (nav);
  \draw[arr] (mission) -| node[lbl, pos=0.4] {odgoda 10~s} (task);
  \draw[arr, dashed] (mission) -| (rviz);

  \draw[arr] (sim) -- (gz);
  \draw[arr] (sim) -- (spwn);
  \draw[arr] (sim) -- (bridge);
  \draw[arr] (sim) -- (ctrl);

  \draw[arr] (nav) -- (amclnode);
  \draw[arr] (nav) -- (navstack);
  \draw[arr] (nav) -- (navhelper);
  \draw[arr] (nav) -- (navsafety);

  \draw[arr] (task) -- (mg);
  \draw[arr] (task) -- (aruco);
  \draw[arr] (task) -- node[font=\tiny, fill=white, inner sep=1pt] {odgoda 3~s} (mtask);
\end{tikzpicture}

  \caption{Hijerarhijski slijed pokretačkih skripti s vremenskim razmacima inicijalizacije}
  \label{fig:launch-hijerarhija}
\end{figure}

U trenutku $t = 0$~s \texttt{sim.launch.py} pokreće simulator Gazebo, učitava URDF model preko
xacro procesora, stvara robota u ishodištu svijeta (početna AMCL poza $(0, 0, 0)$), uspostavlja
\texttt{ros\_gz\_bridge} te pokreće spawner procese za svih osam kontrolera s rokom čekanja od 120~s.
Nakon 6~s \texttt{nav2.launch.py} podiže kartu, AMCL lokalizaciju, Nav2 poslužitelje, potencijalna
polja i sigurnosne monitore. U trenutku $t = 8$~s opcionalno se otvara RViz2, a na $t = 10$~s
\texttt{task.launch.py} aktivira \texttt{move\_group} i detektore markera. Konačno, na $t = 13$~s
pokreće se \texttt{main\_task} s argumentom \texttt{mission:=true}, čime započinje autonomna misija.

\subsection{Struktura repozitorija, ovisnosti i projektno okruženje}
\label{sec:02-repozitorij}

% izvor: ros2.repos:1-35; patches/*; scripts/run_native.sh; README.md §8-9
Izvorni kod projekta podijeljen je u četiri vlastita paketa unutar direktorija \texttt{src/}:
\begin{itemize}
  \item \texttt{pas\_dual\_arm\_bringup}: pokretačke skripte cijelog sustava, konfiguracijske YAML datoteke za kontrolere, Nav2 i mostove, te definicija simulacijskog svijeta i karata.
  \item \texttt{pas\_dual\_arm\_scripts}: implementacija svih misijskih, navigacijskih i percepcijskih čvorova u programskom jeziku Python.
  \item \texttt{pas\_dual\_arm\_description}: cjeloviti xacro modeli sklopljenog robota, omnidirekcijske baze, senzorskih elemenata i kolizijskih geometrija.
  \item \texttt{pas\_dual\_arm\_moveit\_config}: semantički opisi robota (SRDF), konfiguracija kinematičkih rješavača KDL i OMPL algoritama za dvoručni sustav.
  \item \texttt{dual\_arm\_torso}: CAD meshevi i kinematički opis linearnih vodilica torza \cite{mail2026}.
\end{itemize}

% izvor: ros2.repos:20-35; patches/
Vanjski paketi definirani su u datoteci \texttt{ros2.repos} i vezani na točno određene commitove:
\texttt{ros2\_kortex} (\texttt{116d87a}), \texttt{omni\_base\_simulation} (\texttt{77248ac}),
\texttt{pan\_tilt\_ros} (\texttt{9b08758}), \texttt{realsense-ros} (\texttt{6d87b07}) i
\texttt{aruco\_ros} (\texttt{86a0bbb}). Nad vanjskim paketima primijenjene su dvije lokalne
zakrpe u direktoriju \texttt{patches/}: \texttt{pan\_tilt\_ros} zakrpa koja dodaje korektne inercijske
matrice i granice momenta za Gazebo Fortress fiziku te zakrpa za \texttt{ros2\_kortex} koja uklanja
nepodržane Isaac Gym argumente iz makroa Robotiq 2F-85 hvataljke.

% izvor: scripts/run_native.sh; notes/03_problemi/P-01_shell_zenoh_contamination.md
Za pouzdano i determinističko izvođenje razvijena je skripta \texttt{./scripts/run\_native.sh}, koja
postavlja izoliranu okolinu, definira \texttt{RMW\_IMPLEMENTATION=rmw\_cyclonedds\_cpp} i
\texttt{IGN\_GAZEBO\_RESOURCE\_PATH} te sprječava interferenciju sa sustavskim varijablama i vanjskim
middleware protokolima (npr. Zenoh, zabilježen u problemu P-01).
